\chapter{Jira Hata Raporlama Akışı}

\section{JiraBugReportBuilder}
Java tarafında \texttt{JiraBugReportBuilder} fail anında Türkçe Jira raporu üretir. Özet, test adımları, beklenen/gerçekleşen sonuç ve teknik metadata içerir.

\section{Türkçe Rapor Şablonu}
Dashboard Jira modal formu şu alanları sunar:
\begin{itemize}
  \item Proje anahtarı (varsayılan: \texttt{LTPLAY})
  \item Öncelik, önem, kanal, QA süreci
  \item Test adımları, beklenen sonuç, gerçekleşen sonuç
  \item Açıklama ve teknik hata detayı
\end{itemize}

\section{Ekran Görüntüsü ve Log Ekleri}
Fail anında screenshot \texttt{screenshots/} altına kopyalanır. \texttt{FailureLogWriter} teknik log üretir. Dashboard ``Ekleri İndir (.zip)'' ile her iki dosyayı birlikte indirmeyi destekler.

\section{localStorage ve jira-reports.json}
``Hatayı Raporla'' ile kayıt \texttt{qa-dashboard-reported-errors} localStorage anahtarına yazılır. Sunucu tarafı entegrasyon için \texttt{jira-reports.json} opsiyonel olarak kullanılabilir.

% TODO: Örnek Jira rapor metnini Ek C olarak ekle.
